\documentclass[12pt]{article}
\usepackage[utf8]{inputenc}
\usepackage[T1]{fontenc}
\usepackage[portuguese]{babel}
\usepackage[a4paper,margin=3cm]{geometry}
\usepackage{graphicx}
\usepackage{booktabs}
\usepackage{array}
\usepackage{float}
\usepackage{setspace}
\usepackage{siunitx}
\usepackage{hyperref}
\sisetup{output-decimal-marker={,},group-separator={.},group-minimum-digits=3}

\title{Projeto Final de Ciência de Dados para Engenheiros\\Relatório Técnico -- Módulo 2}
\author{Vinícius César Sena Torres -- RA 22409225}
\date{\today}

\begin{document}
\maketitle

\begin{abstract}
Este documento descreve a mineração de regras associativas para identificar padrões de coocorrência entre os produtos mais vendidos do e-commerce. O algoritmo Apriori foi aplicado ao conjunto de \num{5789} cestas (usuários) contendo os 50 itens mais populares, com suporte mínimo de 3\% e confidência mínima de 30\%. As regras extraídas apresentam lift acima de 1,05, assegurando ganho real sobre a compra independente.
\end{abstract}

\section{Introdução}
A identificação de itens comprados em conjunto sustenta promoções cruzadas, kits e recomendações personalizadas. Utilizamos a mesma base do Módulo 1 e nos concentramos nos produtos líderes (P00265242 sozinho aparece em \num{32.5}\% das cestas) para garantir significância estatística e facilidade de implementação comercial.

\section{Metodologia}
\begin{itemize}
    \item \textbf{Seleção de itens}: recortei os 50 produtos com maior frequência; juntos eles representam \num{84}\% de todas as transações do período.
    \item \textbf{Formatação das cestas}: construí matriz binária \(5789 \times 50\), onde cada linha corresponde a um usuário e cada coluna indica a presença do produto.
    \item \textbf{Geração de itemsets}: Apriori com suporte mínimo de 3\% resultou em \num{3028} combinações frequentes.
    \item \textbf{Filtragem de regras}: calculei confidência e lift, retendo \num{6230} regras com confidência \(\geq 30\%\) e lift superior a 1,05.
    \item \textbf{Visualizações}: gráficos de suporte vs. confidência e mapa de calor das doze principais referências para interpretar hubs de coocorrência (familias de produtos 5 e 1).
\end{itemize}

\section{Resultados e Discussão}
A Tabela~\ref{tab:regras} apresenta as cinco regras com maior lift. Todas envolvem pares de produtos que aumentam substancialmente a probabilidade de adquirir itens complementares, em especial o código P00111142.

\begin{table}[H]
    \centering
    \caption{Principais regras associativas}
    \label{tab:regras}
    \begin{tabular}{p{5cm}p{3cm}ccc}
        \toprule
        Antecedentes & Consequente & Suporte (\%) & Confidência (\%) & Lift \\
        \midrule
        P00102642, P00112542 & P00111142 & \num{3.18} & \num{49.9} & \num{2.84} \\
        P00086442, P00112542 & P00111142 & \num{3.13} & \num{49.3} & \num{2.81} \\
        P00058042, P00080342 & P00113242 & \num{3.07} & \num{44.6} & \num{2.77} \\
        P00112542, P00145042 & P00111142 & \num{3.39} & \num{48.5} & \num{2.77} \\
        P00112142, P00112542 & P00111142 & \num{3.78} & \num{48.5} & \num{2.76} \\
        \bottomrule
    \end{tabular}
\end{table}

Observa-se que o produto P00112542 aparece em quatro das cinco regras principais, funcionando como hub de conveniência: sempre que combinado com P00102642, P00145042 ou P00112142, a probabilidade de P00111142 ser comprado aumenta em quase 180\% (lift $>2.7$). Já a dupla (P00058042, P00080342) é altamente preditiva para P00113242, sugerindo kits naturais para campanhas de eletrônicos.

O diagrama suporte vs. confidência mostrou concentração de regras sólidas entre 3\% e 5\% de suporte, faixa considerada aplicável para promoções em massa. O mapa de calor das coocorrências confirmou que itens das categorias 5 e 1 compartilham boa parte dos clientes, reforçando a necessidade de ações coordenadas entre essas famílias.

\section{Recomendações}
\begin{enumerate}
    \item Criar pacotes promocionais envolvendo P00112542 e P00111142, com descontos progressivos ou sugestões automáticas no carrinho, dado o lift superior a 2,8.
    \item Utilizar as regras com suporte \(>3\%\) como gatilhos de recomendações no checkout e em campanhas digitais, pois atingem massa crítica de clientes.
    \item Monitorar trimestralmente o portfólio: novos hubs podem surgir; caso o mix se diversifique, reduzir o suporte mínimo para 2\% garantirá a descoberta de nichos emergentes.
\end{enumerate}

\section{Conclusão}
O módulo de regras associativas estruturou evidências quantitativas para promoções cruzadas. A seleção criteriosa dos itens e o controle de suporte/confidência asseguram que os insights não sejam fruto de ruído e podem ser traduzidos diretamente em banners, recomendações e políticas de precificação dinâmica.

\begin{thebibliography}{9}
\bibitem{dataset} Kaggle. \emph{Black Friday Data Set}. Disponível em: \url{https://www.kaggle.com}. Acesso em: \today.
\bibitem{han} HAN, Jiawei; PEI, Jian; KAMBER, Micheline. \emph{Data Mining: Concepts and Techniques}. 4. ed. Elsevier, 2022.
\end{thebibliography}

\end{document}
